\documentclass[a4paper]{article}

\usepackage{graphicx}
\usepackage{amsmath,amssymb}
\usepackage{minted}
\usepackage{multirow}
\usepackage{float}
\usepackage{todonotes}
\usepackage{forest}
\usepackage{xstring}
\usepackage{xcolor}
\usepackage[most]{tcolorbox}
\usepackage{xparse}
\usepackage{natbib}

\usetikzlibrary{graphs,graphdrawing}
\usegdlibrary{layered}

% for external graphviz
\input{/home/donkarlo/Dropbox/repo/nd_ai_project/src/nd_ai/knoweldge_representation/language/formal/computer/programming/tex/latex/code_clips/external_graphviz.tex}

% for tags
\input{/home/donkarlo/Dropbox/repo/nd_ai_project/src/nd_ai/knoweldge_representation/language/formal/computer/programming/tex/latex/parts/control_sequence/kind/semtag/semtag.tex}

% for links
\input{/home/donkarlo/Dropbox/repo/nd_ai_project/src/nd_ai/knoweldge_representation/language/formal/computer/programming/tex/latex/code_clips/seehere.tex}

\title{}
\date{}

\begin{document}
    \maketitle
    
    \cite{gordon-1993-novel-approach-to-nonlinear-non-gaussian-bayesian-state-estimation} introduced particle filter as a sequential Monte Carlo method used to estimate the hidden state of a dynamic system from noisy observations.
    
    Assume that the hidden state at time $t$ is $x_t$ and the observation is $y_t$.
    
    \begin{align}
        x_t &\sim p(x_t \mid x_{t-1}) \\
        y_t &\sim p(y_t \mid x_t)
    \end{align}
    
    The first equation describes how the state changes over time.
    The second equation describes how the observation is generated from the hidden state.
    
    Instead of representing the posterior distribution exactly,
    
    \begin{equation}
        p(x_t \mid y_{1:t})
    \end{equation}
    
    a particle filter approximates it using many weighted samples, called particles:
    
    \begin{equation}
        p(x_t \mid y_{1:t}) \approx \sum_{i=1}^{N} w_t^{(i)} \delta(x_t - x_t^{(i)})
    \end{equation}
    
    where $x_t^{(i)}$ is the $i$-th particle, $w_t^{(i)}$ is its weight, and $N$ is the number of particles.
    
    At each time step, the algorithm has three main steps.
    
    
    \section{Prediction}
        
        Each particle is moved forward using the state transition model:
        
        \begin{equation}
            x_t^{(i)} \sim p(x_t \mid x_{t-1}^{(i)})
        \end{equation}
    
    
    \section{Update}
        
        Each particle receives a weight according to how well it explains the new observation:
        
        \begin{equation}
            w_t^{(i)} \propto p(y_t \mid x_t^{(i)})
        \end{equation}
        
        Then the weights are normalized:
        
        \begin{equation}
            w_t^{(i)} =
            \frac{w_t^{(i)}}{\sum_{j=1}^{N} w_t^{(j)}}
        \end{equation}
    
    
    \section{Resampling}
        
        Particles with large weights are copied more often, and particles with small weights are removed.
        
        \begin{equation}
            \{x_t^{(i)}, w_t^{(i)}\}_{i=1}^{N}
            \rightarrow
            \{\tilde{x}_t^{(i)}, \frac{1}{N}\}_{i=1}^{N}
        \end{equation}
        
        After resampling, all particles have equal weight again.
    
    
    \section{Simple Interpretation}
        
        A particle filter represents belief about the hidden state by many guesses.
        Each guess is tested against the new observation.
        Good guesses survive, and bad guesses disappear.
        Over time, the cloud of particles follows the probable state of the system.
        
        \bibliographystyle{plainnat}
        \bibliography{/home/donkarlo/Dropbox/repo/nd_shared_research/bibliography/bibliography.bib}
\end{document}